% !TeX root = ./../main.tex

\chapter{Podsumowanie}\label{chapter:summary}

\agentinfo{
Rozdział pisać na końcu.
Każde twierdzenie musi mieć pokrycie w wynikach rozdziałów 4--8; nie wprowadzać tu nowych wyników.
}

\section{Najważniejsze wyniki}\label{sec:summary:results}
\agentinfo{Po jednym akapicie na rozdział badawczy, z kluczowymi wartościami liczbowymi.}

\section{Weryfikacja tez badawczych}\label{sec:summary:theses}
\agentinfo{
Dla każdej tezy z wprowadzenia: stwierdzenie, czy została potwierdzona, w jakim zakresie i na podstawie których wyników.
Odpowiedzi na pytania badawcze PB1--PB6.
}

\section{Ograniczenia}\label{sec:summary:limitations}
\agentinfo{
Zebrać ograniczenia z dyskusji poszczególnych rozdziałów: skala eksperymentów, wycofane urządzenia, zależność od kalibracji, brak eksperymentów na sprzęcie fotonicznym, rozmiar zbiorów danych.
}

\section{Kierunki dalszych badań}\label{sec:summary:future}
\agentinfo{
Kierunki zapowiadane w artykułach: implementacja na układach fotonicznych, komunikacja kwantowa, przetwarzanie obrazu w stanie zakodowanym, wykrywanie obiektów (przyrządy lotnicze), rozszerzenie modelu szumu na układy wielokubitowe, wbudowanie diagnostyki łamania symetrii w schematy odporne na błędy.
}
